\subsection{凸锥的极生成元}

设 C 是向量空间 E 中以 0 为顶点的一个凸锥；显然 C 中除顶点外没有别的点能是极点；顶点是 C 的极点，当且仅当 C 是真的且尖的。

定义 2. --- 设 C 是向量空间 E 中以 0 为顶点的一个凸锥。若每条含于 C、不含 0 且与 D 相交的开线段都含于 D，则称一条从 0 出发的半线 D ⊂ C 为 C 的一个极生成元。

这等价于说：对所有 \(x \in D\)（满足 \(x \neq 0\)），若 \(y \neq 0\)、\(y' \neq 0\) 是 C 中满足 \(x = y + y'\) 的两个点，则必有 \(y \in D\) 与 \(y' \in D\)。

注 1. --- 设 C 是 E 中的一个真尖凸锥，并在 E 上考虑使 C 成为 \(\geqslant 0\) 元素集的序结构（II, p.~12, 命题 13）；为使 E 的一个元素——设为 x \textgreater{} 0——属于 C 的某个极生成元，其必要充分条件是：每个以 x 为上界的元素 \(y \geqslant 0\) 都具有形式 \(\lambda x\)，其中 \(0 \leqslant \lambda \leqslant 1\)：事实上，说 y 以 x 为上界，意味着 \(x = y + y'\)，其中 \(y' \in C\)，由此即得结论。

命题 3. --- 在向量空间 E 中，设 C 是以 0 为顶点的凸锥，\(x_{0} \neq 0\) 是 C 的一点，D 是含于 C、从 0 出发且包含 \(x_{0}\) 的一条半线。设 H 是包含 \(x_{0}\) 且不通过 0 的一个超平面。则 D 是 C 的极生成元，当且仅当 \(x_{0}\) 是 H ∩ C 的一个极点。

该条件显然是必要的。反之，设它成立；假设存在一条直线 \(D'\)，它不含 D、通过 \(x_{0}\)，且使 \(D' \cap C\) 包含一个含 \(x_{0}\) 的开线段。设 \(y \neq 0\) 是 \(D'\) 的一个方向向量；诸假设蕴含：点 \((1 + \lambda)x_{0} + \mu y\) 当 \(|\lambda|\) 与 \(|\mu|\) 充分小时属于 C。但在由 D 与 \(D'\) 所决定且赋有典范拓扑的平面 P 中，\(x_{0}\) 是 \(P \cap C\) 的一个内点，由此推出，直线 \(P \cap H\) 包含一个含于 \(H \cap C\) 且含 \(x_{0}\) 的开线段。这与假设矛盾。

定义 3. --- 设 C 是 Hausdorff 拓扑向量空间 E 中的一个凸集。C 的一个非空紧凸集 A 称为 C 的一个帽，如果 A 在 C 中的补集 C−A 是凸的。

设 C 是 E 中以 0 为顶点的一个尖凸锥，并设 A 是 C 的一个帽。写 B = C - A。对每条从 0 出发的闭半线 L ⊂ C，集 L ∩ A 与 L ∩ B 都是凸集，它们在 L 中互为补集，其并为 L，且 L ∩ A 是紧的。由于 L ∩ A 至少对一条半线 L 非空，我们看出 0 ∈ A，从而 L ∩ A 是以 0 为一个端点的闭线段。若 A 存在，则 C 是真锥。

命题 4. --- 设 C 是 E（一个 Hausdorff 局部凸空间）中以 0 为顶点的一个尖凸锥。

\begin{enumerate}
\def\labelenumi{\alph{enumi})}
\tightlist
\item
  设 A 是 C 的一个帽。设 p 是 A 的度规到 C 上的限制（II, p.~20）。\(x \in C\) 中使 \(p(x) \leqslant 1\) 者所成之集就是集 A。函数 p 是下半连续的，且具有下列性质：
\end{enumerate}

\begin{enumerate}
\def\labelenumi{(\roman{enumi})}
\item
  对任意的 \(x, y\)（在 \(\mathbf{C}\) 中），我们有 \(p(x + y) = p(x) + p(y)\)。
\item
  对任意的 \(x \in \mathbf{C}\) 与 \(\lambda \in \mathbf{R}_{+}^{*}\)，我们有 \(p(\lambda x) = \lambda p(x)\)。
\item
  若 \(x \in C\)，则关系 \(p(x) = 0\) 等价于 x = 0。
\end{enumerate}

\begin{enumerate}
\def\labelenumi{\alph{enumi})}
\setcounter{enumi}{1}
\tightlist
\item
  反之，设 \(p\) 是定义在 \(\mathbf{C}\) 上、取值于 \([0, +\infty]\) 的函数，满足 a) 的条件 (i)、(ii)。设 \(\mathbf{A}\) 为由 \(x \in \mathbf{C}\) 中使 \(p(x) \leqslant 1\) 者组成的集。则 \(\mathbf{A}\) 与 \(\mathbf{C} - \mathbf{A}\) 都是凸的。A 是帽，当且仅当 \(\mathbf{A}\) 紧且非空。
\end{enumerate}

陈述 \(b)\) 是明显的。\(a)\) 中所述的性质是命题 4 前面的那些注记以及 II, p.~20 的命题 22 与 II, p.~20 的命题 23 的推论，但下式除外：

\[
p (x + y) \geqslant p (x) + p (y).
\]

只需对 \(x \neq 0\) 且 \(y \neq 0\) 的情形证明最后这个不等式；于是有 \(p(x) > 0\)、\(p(y) > 0\)。设 \(\mu, \lambda\) 是两个 \textgreater{} 0 的数，满足 \(\lambda < p(x)\)、\(\mu < p(y)\)，并用 B 表示 A 在 C 中的补集。我们有 \(x \in \lambda B\)、\(y \in \mu B\)，因此 \(x + y \in \lambda B + \mu B\)；由 B 的凸性，有 \(\lambda B + \mu B \subset (\lambda + \mu)B\)，从而 \(p(x + y) > \lambda + \mu\)，这就蕴含上面所述的不等式。

推论 1. --- 设 C 是 E（一个 Hausdorff 局部凸空间）中以 0 为顶点的尖凸锥，并设 p 是 C 的一个帽 A 的度规。则 A 的极点就是点 0，以及 C 的诸极生成元上满足 \(p(x) = 1\) 的那些点 x。

显然 0 是 A 的极点。设 x 是 C 的一个极生成元 L 上的点，且满足 \(p(x) = 1\)。设 y、z 是 A 中满足 \(x = \frac{1}{2}(y + z)\) 的两点。由于 L 是极的，我们有 \(y = \lambda x\) 与 \(z = \mu x\)，其中 \(\lambda\) 与 \(\mu\) 是 \(\geqslant 0\) 的数，满足 \(\frac{1}{2}(\lambda + \mu) = 1\)、\(\lambda = \lambda p(x) = p(y) \leqslant 1\) 与 \(\mu = \mu p(x) = p(z) \leqslant 1\)，由此得 \(\lambda = \mu = 1\)，从而 y = z = x；所以 x 是 A 的极点。反之，设 \(x \neq 0\) 是 A 的一个极点。显然 \(p(x) = 1\)。设 y、\(y'\) 是 C 中满足 \(x = y + y'\) 的两点，我们将证明 y、\(y'\) 与 x 成比例。不失一般性，可以假设数 \(\lambda = p(y)\) 与 \(\lambda' = p(y')\) 有限且 \textgreater0。于是由命题 4 (i)，\(\lambda^{-1}y \in A\)、\(\lambda'^{-1}y' \in A\)、\(\lambda + \lambda' = 1\)，且等式 \(x = \lambda(\lambda^{-1}y) + \lambda'(\lambda'^{-1}y')\) 按假设蕴含

\[
x = \lambda^ {- 1} y = \lambda^ {\prime - 1} y ^ {\prime}.
\]

推论 2. --- C 中属于 C 的某个帽的每一点，也属于 C 的诸极生成元之并的闭凸包。

这由推论 1 与 Krein-Milman 定理（II, p.~55, 定理 1）立得。

* 例. --- 设 X 是一个局部紧且 \(\sigma\) -紧的空间。设 C 是 \(\mathcal{M}_{+}(X)\)（赋淡拓扑）中以 0 为顶点的一个闭凸锥。我们将证明 C 是它的诸帽之并。设 \((X_{n})\) 是 X 的开相对紧集的一个递增序列，其并为 X。设 \(\mu\) 是 C 中 \(\neq 0\) 的一个元素。存在 \(\alpha_{n} > 0\) 使 \(\sum \alpha_{n}\mu(X_{n}) = 1\)。

对每个测度 \(v \in C\)，令 \(p(v) = \sum_{n} \alpha_{n} v(X_{n}) \in [0, +\infty]\)。C 上的函数 p 满足命题 4 的条件 (i) 与 (ii)。它对淡拓扑是下半连续的（INT, IV, 第 2 版, § 1, No.~1, 命题 4）。于是 \(\gamma \in C\) 中使 \(p(\gamma) \leqslant 1\) 者所成的集 A 是闭的且非空。另一方面，X 的每个紧集都含于某个 \(X_{n}\) 之中，于是 A 既为淡有界，从而也是淡紧的（INT, III, 第 2 版, § 1, No.~9, 命题 15）。因此集 A 是 C 的一个包含 \(\mu\) 的帽。

命题 5. --- 设 C 是 E（一个 Hausdorff 弱空间）中以 0 为顶点的一个真凸锥；假设 C 对由 E 的一致结构所诱导的一致结构是完备的，并且 C 中 0 有一个可数的基本邻域系。则 C 是它的诸帽之并，并且是其诸极生成元之并的闭凸包。

第二个断言由第一个断言及上面的推论 2 得出。利用 II, p.~52 的命题 11，可把命题化归为 \(\mathbf{E} = \mathbf{R}^{\mathrm{I}}\) 且 \(\mathbf{C} \subset \mathbf{R}_{+}^{\mathrm{I}}\) 的情形。对所有 \(\alpha \in \mathbf{I}\)，把投影 \(\operatorname{pr}_{\alpha}\)（\(\mathbf{E}\) 中的）记作 \(f_{\alpha}\)；于是 \(f_{\alpha}\) 是连续线性型。另一方面，设 \((\mathrm{V}_n)_{n \in \mathbb{N}}\) 是 \(\mathbb{C}\) 中 0 的一个可数基本邻域系。按 \(\mathbf{E}\) 的拓扑的定义，对每个 \(n \in \mathbb{N}\)，存在一个有限子集 \(\mathrm{J}_n\)（\(\mathbf{I}\) 的）与一个数 \(\varepsilon_n > 0\)，使得 \(\mathrm{V}_n\) 包含 \(\mathrm{W}_n\)——即 \(x \in \mathbb{C}\) 中使 \(f_{\alpha}(x) \leqslant \varepsilon_n\)（对所有 \(\alpha \in \mathrm{J}_n\)）者所成的集；令 \(\mathrm{J} = \bigcup_{n \in \mathbb{N}} \mathrm{J}_n\)。设 \(y \neq 0\) 是 \(\mathbb{C}\) 中的一点，并设 \(p\) 是函数 \(\sum_{\alpha \in \mathbf{J}} \lambda_{\alpha}(f_{\alpha}|\mathbb{C})\)，其中诸 \(\lambda_{\alpha} > 0\) 取得使 \(p(y) = 1\)；这是可能的，因为若 \(f_{\alpha}(y) = 0\) 对所有 \(\alpha \in \mathbf{J}\) 成立，则 \(y \in \mathrm{V}_n\) 对所有 \(n\) 成立，这就蕴含 \(y = 0\)，与假设相违。现在我们注意到，对所有 \(\alpha \in \mathbf{I}\)，函数 \(f_{\alpha}|\mathbb{C}\) 在点 0 处连续，因此存在一个 \(n \in \mathbb{N}\)，使 \(f_{\alpha}\) 在某个 \(\mathrm{W}_n\) 中有界，从而在 \(\mathbb{C}\) 中以有限多个函数 \(f_{\beta}|\mathbb{C}\)（其中 \(\beta \in \mathbf{J}\)）的一个线性组合为其一个上界。由此推出，若 \(\mathbf{A}\) 是由 \(x \in \mathbb{C}\) 中使 \(p(x) \leqslant 1\) 者组成的集，则 \(f_{\alpha}\) 在 \(\mathbf{A}\) 中有界（对所有 \(\alpha \in \mathbf{I}\)）。由于 \(p\) 在 \(\mathbb{C}\) 中下半连续，可知 \(\mathbf{A}\) 在 \(\mathbb{C}\) 中闭且非空，从而是紧的。既然显然 \(p\) 满足 II, p.~58 命题 4 的条件 (i) 与 (ii)，我们就看出 \(\mathbf{A}\) 是 \(\mathbb{C}\) 中的一个帽，且包含 \(y\)。

注 2. --- 存在弱完备且没有极生成元的真凸锥（II, p.~92, 习题 31）。

